\setcounter{theorem}{19}
\section{Hardness of Generators for torus actions}\label{sec:torus}

Let $\C^*$ denote the multiplicative group consisting of all non-zero
complex numbers. A direct product $\T_n = (\C^*)^n$ is called a torus, and
is clearly an abelian group. Tori are important examples of reductive
groups -- any abelian connected reductive group is a torus! It is often the case that it is easier to understand tori in comparison with more general (non-abelian) reductive groups. This is no different for invariant theory, see for example \cite{DK,Wehlau}. We also point to \cite[Proposition~3.3]{DM-exp} for an elementary linear algebraic description of the invariant ring for torus actions. Conjecture 6 of the original publication already fails in this well behaved setting.
This is the content of our Theorem~\ref{thm:torus-action}, which we will prove in this section.
Recall that $\ST_n(\C) \cong \{ \bz \in \T_n : z_1\cdots z_n = 1 \}$, which is itself a torus.

\begin{theorem}[Theorem~\ref{thm:torus-action}, restated]
	Consider the natural action of $G = \ST_n(\C) \times \ST_n(\C) \times \ST_n(\C)$
	on $V = \C^n \otimes \C^n \otimes \C^n$, where an element
	$(\ba, \bb, \bc) \in G$ acts on a tensor $u \in V$ as follows:
  $(\ba,\bb,\bc) \cdot u := v$, such that
	$v_{ijk} = a_i b_j c_k u_{ijk}$. Any set of generators
	for the invariant ring of this action cannot have a polynomial sized
	(in $n$) succinct encoding, unless ${\rm NP \subseteq P/poly}$.
\end{theorem}

\begin{proof}
	Suppose that the natural action above has a set of generators with
	a polynomial sized succinct encoding. Thus, there is an arithmetic
	circuit $\cC(\bx, \by)$ of size $s = {\rm poly}(n)$, where
	$\bx = (x_{ijk})_{i,j,k=1}^n$ is the set of variables corresponding
	to $V$ and $\by = (y_1, \ldots, y_r)$ is the set of auxiliary
	variables, with $r = {\rm poly}(n)$.
	Moreover, from the definition of the size of a succinct encoding we also have that the
	constants used in the computation of $\cC(\bx, \by)$ are rational numbers with bit
	complexity bounded by $b = {\rm poly}(n)$. In particular, $\cC(\bx, \by) \in \Q[\bx,\by]$.

	Let us consider the circuit $\cC(\bx, \by)$ as a circuit whose constants
	are in $\Q[\by]$ and whose variables are only the $\bx$ variables,
	that is, a circuit in $\Q[\by][\bx]$.
  Then, Proposition~\ref{prop:hom-components}
	tells us that there exists a homogeneous circuit $\cC_n(\bx, \by)$,
	in the $\bx$ variables, of degree $n$ and size $O(n^2 s)$ that computes the homogeneous component of $\cC(\bx,\by)$ of degree $n$ as a function of~$\bx$.
  Moreover, the constants of this circuit are a subset of the constants used in the circuit $\cC(\bx, \by)$.
  Since we consider the latter as a circuit in only the $\bx$ variables, the constants in this case are given by the elements of $\Q$ used in the computation of
  $\cC$ as well as the auxiliary variables $\by$.
  In particular, $\cC_n(\bx, \by)$ can be written in the following way:
	\begin{equation}\label{eq:deg-n-cks-1}
		\cC_n(\bx, \by) = \sum_{\bbm \in \cN_n(\bx)} f_\bbm(\by) \cdot \bbm,
	\end{equation}
	where $\cN_n(\bx)$ is the set of all monomials of degree $n$ in the
	variables $\bx$ and $f_\bbm(\by)$ are polynomials in the variables $\by$ of
	degree at most~$2^s$, as the circuit $\cC$ has size at most $s$.

	In Proposition~\ref{prop:least-degree-torus} below, we will show that the invariants
	of minimum degree of our action are in degree $n$, and these
	are spanned by the (maximum) $3$-dimensional matching monomials.
  Thus, if a monomial of degree $n$ is invariant under our action, it must be the case that this monomial corresponds to a $3$-dimensional matching.
  Moreover, the action maps any monomial (invariant or not) to a constant times itself. 
  As $\cC_n(\bx, \by)$ must only compute invariant polynomials, this
	implies that equation~\eqref{eq:deg-n-cks-1} is actually of the following form:
	\begin{equation}\label{eq:deg-n-cks-2}
		\cC_n(\bx, \by) = \sum_{\bbm \in \cM_n(\bx)}
		f_\bbm(\by) \cdot \bbm,
	\end{equation}
	where $\cM_n(\bx)$ is the set of all $3$-dimensional matching monomials
	in the variables $\bx$.
  Moverover, since $\cC(\bx,\by)$ succinctly encodes of a set of generators, the span of $\{\cC_n(\bx,\alpha)\}_{\alpha\in\C^r}$ must necessarily be the same as the span of the 3-dimensional matching polynomials.

	We will now show that the existence of the circuit $\cC_n(\bx, \by)$ implies that
	${\rm NP\subseteq P/poly}$. For that purpose, we will show that given $\cC_n(\bx, \by)$ one can
	solve the $3$-dimensional matching problem in ${\rm P/poly}$.
	Let $H$ be a tripartite 3-uniform hypergraph, whose edges are given by a subset
	$E \subseteq [n] \times [n] \times [n]$.
	We can associate to this graph the tensor $v \in V$ where
	$v_{ijk} = 1$ if hyperedge $(i,j,k) \in E$ and $v_{ijk} = 0$ otherwise.
	Note that $H$ has a $3$-dimensional matching of size $n$ if and only if at least one
	of the $3$-dimensional matching monomials \emph{does not vanish} on our tensor $v$.
	This last
	condition is equivalent to the fact that the circuit $\cC_n(v, \by)$ does not compute the
	zero polynomial (as we know that the span of the set $\{\cC_n(\bx, \alpha)\}_{\alpha \in \C^r}$
	is the same as the span of all $3$-dimensional matching monomials).
	Thus, to solve the $3$-dimensional
	matching problem in ${\rm P/poly}$ it is enough to give a randomized polynomial time algorithm
	for testing whether $\cC_n(v, \by)$ is the zero polynomial or not.\footnote{It is enough to
	give a randomized polynomial time algorithm because we know that ${\rm BPP/poly = P/poly}$.}

	Since $\cC_n(v, \by)$ is a circuit of size ${\rm poly}(n)$ with rational constants of bit
	complexity ${\rm poly}(n)$, it computes a polynomial $P(\by)$ with rational coefficients
	having bit complexity at most $\exp(\poly(n))$ and degree at most $\exp(\poly(n))$.
  This is the setting in which
	Theorem~\ref{thm:PIT-BPP} applies, giving us the desired randomized polynomial
	time algorithm. This concludes our proof modulo Proposition~\ref{prop:least-degree-torus},
	which we will now turn our attention to.
\end{proof}



%
%
%
%

%
%
%
%
%
%

%
%
%
%
%
%
%
%

\noindent
In the following proposition, we denote by $\mathcal{S}_n$ the symmetric group on $n$ letters.


\begin{proposition}\label{prop:least-degree-torus}
	The maximum $3$-dimensional matching monomials  $\prod_{i=1}^n x_{i \sigma(i) \tau(i)}$,
	where $\sigma, \tau \in \mathcal{S}_n$, span the invariants of degree
	$n$ of the natural action of $G = \ST_n(\C) \times \ST_n(\C) \times \ST_n(\C)$
	on $V = \C^n \otimes \C^n \otimes \C^n$. Moreover, there are no nonconstant invariants of degree less than $n$ for this action.
\end{proposition}

\begin{proof}
  Since the action maps any monomial to a constant times itself, it is easy to see that the invariant polynomials are generated by invariant monomials.
	To prove the proposition, it is therefore enough to show that the matching monomials are invariant, that there are no other invariant monomials of degree~$n$, and that there are no invariant monomials of smaller degree.

  We first prove that the matching monomials are invariant.
	Note that the natural action of $G$ on $V$ induces the following action on the variables $x_{ijk}$:
	$(\ba, \bb, \bc) \cdot x_{ijk} = (a_i b_j c_k)^{-1} \cdot x_{ijk}$.
  \footnote{The inverse comes from the general formula $(g\cdot p)(v) := p(g^{-1} \cdot v)$. }
  Additionally, note that
	$\prod_{\ell=1}^n a_\ell = \prod_{\ell=1}^n b_\ell = \prod_{\ell=1}^n c_\ell = 1$. Given a
	matching monomial $\prod_{i =1}^n x_{i \sigma(i) \tau(i)}$, we therefore have that
	\begin{align*}
		(\ba, \bb, \bc) \cdot \prod_{i =1}^n x_{i \sigma(i) \tau(i)}
		& = \prod_{i =1}^n \left( (a_i b_{\sigma(i)} c_{\tau(i))})^{-1} \cdot x_{i \sigma(i) \tau(i)}  \right) \\
		& = \prod_{i=1}^n (a_i b_{\sigma(i)} c_{\tau(i)})^{-1}  \cdot
		\prod_{i =1}^n x_{i \sigma(i) \tau(i)} \\
		& = \prod_{i =1}^n x_{i \sigma(i) \tau(i)}
	\end{align*}
%
	where in the last equality we note that for any permutation $\sigma \in \mathcal{S}_n$ (or $\tau$)
	we have $1= \prod_{\ell=1}^n a_\ell = \prod_{\ell=1}^n a_{\sigma(\ell)}$ (and similarly for
	$\bb$ and $\bc$). This proves that the matching monomials are invariant monomials of the natural
	$G$-action on $V$.

	Now, let us prove that no other monomial of degree $n$ is an invariant for this action.
	Let $\prod_{m=1}^n x_{i_m j_m k_m}$ be a monomial, where $(i_m, j_m, k_m) \in [n]^3$, that is not a matching monomial.
	Then there exists some coordinate, w.l.o.g.\ the first coordinate, for which the set~$\{i_m\}_{m=1}^n$ is a strict subset of~$[n]$.
  Equivalently, there is an element $\ell \in [n]$ such that $\ell \not\in \{i_m\}_{m=1}^n$.
  W.l.o.g., we can assume that
	$\ell = 1$. Thus, the action of $\ba = (\alpha^{n-1}, \alpha^{-1}, \ldots, \alpha^{-1}),
	\bb = \bc = (1, \ldots, 1)$  on our monomial $\prod_{m=1}^n x_{i_m j_m k_m}$ is as follows:
	\begin{equation*}
		(\ba, \bb, \bc) \cdot \prod_{m=1}^n x_{i_m j_m k_m}
		= \alpha^n \cdot \prod_{m=1}^n x_{i_m j_m k_m}
	\end{equation*}
	which proves that this monomial is not an invariant.
  This completes the proof that the matching
	monomials span the invariants of degree $n$.

	Now we are left with proving that there are no nonconstant monomials of degree less than~$n$ that are invariant.
  Note that if we have a monomial with degree less than $n$, we can represent it as
	$\prod_{m=1}^d x_{i_m j_m k_m}$, where $d < n$ and by the pigeonhole principle, we know that
	there exists $\ell \in [n]$ such that $\ell$ does not appear as a first coordinate entry in
	the set of tuples $\{(i_m, j_m, k_m)\}$. If $d>0$ then, analogously to the previous paragraph, we know
	that such monomials cannot be invariants of the natural action of $G$ over $V$, therefore showing
	that no nonconstant monomial of degree $< n$ can be an invariant. This completes the proof.
\end{proof}





